\chapter{Methodology}\label{chapter:methodology}

\section{Overall Pipeline}

The system follows a linear pipeline from configuration to artefact:

\begin{center}
  Configuration $\!\to\!$ Dataset $\!\to\!$ Backbone $\!\to\!$ Algorithm
  $\!\to\!$ Trainer $\!\to\!$ Checkpoint $\!\to\!$ Sampler $\!\to\!$ Evaluator
  $\!\to\!$ JSONL / Browser UI
\end{center}

Algorithm files own the mathematical objective; all shared modules
(backbone, trainer, sampler, evaluator) remain unmodified across algorithms.
This separation is enforced: algorithm keyword arguments cannot shadow shared
batch, epoch, seed, device, AMP, dataset, backbone, optimizer, evaluation, or
output settings.

\section{Data Preparation}

CIFAR-10 images are resized to $32\!\times\!32$, converted to tensors, and
normalised with mean $(0.5, 0.5, 0.5)$ and standard deviation $(0.5, 0.5, 0.5)$,
mapping pixel values to approximately $[-1, 1]$. Generated tensors are clamped
to $[-1, 1]$ before visualisation and before conversion to uint8 for metric
evaluation.

\begin{figure}[htbp]\centering
  \includegraphics[width=.74\textwidth]{cifar10_examples.png}
  \caption{Representative CIFAR-10 images spanning all ten classes.}\label{fig:cifar}
\end{figure}

The training dataloader shuffles, drops the final incomplete batch, and uses the
configured batch size. The test loader provides real images for FID reference
computation.

\section{Shared Backbone Architecture}

All three algorithms use the same \texttt{SimpleUNet}. Each run instantiates it
from an identical backbone configuration. The network maps:
\begin{equation}
  f_\theta: \mathbb{R}^{3\times 32\times 32} \times [0,1]
    \;\longrightarrow\; \mathbb{R}^{3\times 32\times 32}.
\end{equation}
The scalar time $t$ is encoded by 256-dimensional sinusoidal features and
projected into every residual block via a learned linear layer. The encoder uses
base width 64, channel multipliers $[1, 2, 2]$, and two residual blocks per
stage, with strided convolution downsampling. The decoder mirrors the encoder
with bilinear upsampling and skip connections. Group normalisation (8 groups) and
SiLU activation are used throughout.

\begin{figure}[htbp]\centering
  \includegraphics[width=\textwidth]{architecture.png}
  \caption{Shared SimpleUNet backbone with 6{,}352{,}899 parameters. MF adds a
    separate 323-parameter $r$-embedding that is summed channel-wise with the
    input tensor before the first encoder stage.}\label{fig:arch}
\end{figure}

For MF, the second time endpoint $r$ is mapped through a two-layer projection
with dimensions $1\!\to\!64\!\to\!3$ and a SiLU activation. Its output is added to the input tensor $z$
before the shared backbone. This constitutes 323 additional parameters
(approximately 0.005\% of the backbone).

\section{Training Objectives}

\subsection{Flow Matching}

Given a minibatch of data samples $x_0$ and noise $\varepsilon \sim \mathcal{N}(0, I)$,
a time $t \sim \mathcal{U}(0, 1)$ is sampled, the interpolant $x_t = (1-t)x_0 + t\varepsilon$
is formed, and the loss is:
\begin{equation}
  \mathcal{L}_{\mathrm{FM}}
    = \mathbb{E}\bigl\|v_\theta(x_t, t) - (\varepsilon - x_0)\bigr\|_2^2.
\end{equation}

\subsection{Flow Matching with Logit-Normal Time}

The path, target, and loss formula are identical to FM. The only change is the
time sampling distribution: $t = \sigma(u)$ where $u \sim \mathcal{N}(0, 1)$.
This concentrates training samples near $t=0.5$ and reduces the proportion of
very clean ($t\approx 0$) or very noisy ($t\approx 1$) samples in each batch.

\subsection{Mean Flow}

Given $z_t \sim \mathcal{N}(0, I)$ (start) and $r \sim \mathcal{U}[0, t]$ (interval
start), the network predicts $u_\theta(z_t, r, t)$. The training target is
computed via the Mean Flow Identity:
\begin{equation}
  \hat{u}(z_t, r, t)
    = \underbrace{v_\theta(z_t, t)}_{\text{instantaneous}} \;-\;
      (t - r)\,\underbrace{\frac{d}{dt}\,u_\theta(z_t, r, t)}_{\text{JVP}},
\end{equation}
where the total time derivative is evaluated using \texttt{torch.func.jvp} with
primals $(z_t, r, t)$ and tangents $(v_\theta(z_t, t),\, 0,\, 1)$. The
resulting target is detached from the computational graph before the MSE loss.
With probability $p_\text{same} = 0.25$, the training sets $r = t$, which
recovers the instantaneous velocity loss and stabilises early training.

\section{Optimisation and Checkpointing}

All runs use the AdamW optimiser~\cite{loshchilov2019adamw} with zero weight
decay and no learning rate scheduler. Automatic Mixed Precision (AMP) is enabled
for all runs; MF disables autocast inside the JVP computation to avoid
mixed-dtype derivative errors.

Checkpoints are saved every ten epochs. Each submitted archive preserves the
algorithm module state, optimiser state, configuration, epoch number, and random
seed. MF stores both the backbone and the $r$-embedding state dictionaries. The
archived checkpoints are the model states used for sampling and evaluation.

\section{Sampling and Evaluation}\label{sec:sampling}

After training, each model generates 1{,}000 images at NFE $\in \{1, 5, 10, 20\}$.

\paragraph{FM and FM-LN.}
Sampling starts at $t=1$ (pure noise) and applies backward Euler steps:
\begin{equation}
  x_{t - h} = x_t - h\,v_\theta(x_t, t), \qquad h = 1/\mathrm{NFE}.
\end{equation}

\paragraph{MF.}
The interval $[0, 1]$ is partitioned into NFE equal sub-intervals. For each
sub-interval $[r, t]$, one network call produces the displacement:
\begin{equation}
  z_r = z_t - (t - r)\,u_\theta(z_t, r, t).
\end{equation}
At NFE$=1$, $r=0$ and $t=1$, giving the direct noise-to-image jump.

Wall-clock sampling time is measured with CUDA synchronisation.
TorchMetrics computes FID and IS from images in batches of 32, converting
tensors from $[-1, 1]$ to uint8 in $[0, 255]$ before passing them to the
Inception network. Real-image reference statistics for FID are cached once from
1{,}000 test images alongside a metadata file recording image count, resolution,
and channels; a later evaluation run that requests different settings raises an
explicit error rather than silently reusing a mismatched cache.

\section{Audit Protocol}\label{sec:audit}

Training restarts append new JSONL rows for already-logged epochs. The canonical
rule resolves duplicates by retaining the \textbf{last} occurrence of each
(record\_type, epoch, NFE) key without modifying source files. An audit script
verifies: valid JSON on every row; unique and contiguous canonical epochs;
strictly increasing step and sample counters; finite loss values; constant
backbone parameter count across epochs; and complete NFE coverage in evaluation
records. All three submitted runs pass all checks.

\begin{itemize}
  \item \textbf{FM:} 110 training rows, 12 evaluation rows, 4 sampling rows.
    Duplicate rows cover epochs 1--10 (one restart). Canonical: 100 unique epochs.
  \item \textbf{FM-LN:} 100 training rows, 12 evaluation rows, 4 sampling rows.
    No duplicate training epochs.
  \item \textbf{MF:} 59 training rows, 4 evaluation rows, 4 sampling rows.
    Duplicates cover epochs 1--27 (two restart sequences including an abandoned
    12.5~GB attempt). Canonical: 30 unique epochs, epochs 1--30.
\end{itemize}
